\documentclass[11pt,a4paper]{article}
\usepackage{../../common/elkdoc}

\begin{document}

\labheader{Lab 2 — Logstash}{3.5 h · autonomous · graded deliverable}{Session 4}

\section*{Rules of the game}

You work autonomously, alone or in pairs. Instructors clarify \emph{missions},
never \emph{solutions}. Allowed help: the \textbf{Course 2 slides} and the
\textbf{official Logstash documentation} — nothing else.

\code{scripts/check.sh} is your referee: it prints PASS/FAIL per mission
against the live cluster, and you can run it as often as you like. Everything
you need was covered in Course 2 — when stuck, ask yourself \emph{which slide
was this?}, then find the matching documentation page.

\begin{deliverbox}[Deliverable (graded /20, deadline: \underline{\hspace{2cm}})]
Your final \code{pipeline/pipeline.conf} + a 10-line write-up containing:
(a) the mission 1 explanation, (b) grok vs.\ dissect — what you chose for the
app logs and \emph{why}, (c) your dead-letter strategy and how you verified
that \textbf{nothing is silently lost}.
\end{deliverbox}

\section*{Starter kit \& setup}

\begin{lstlisting}[language=bash]
lab2/
|-- docker-compose.yml         # ES + Kibana + Logstash (pipeline/ auto-reloads)
|-- generator/generate_logs.py # writes the three log sources into logs/
|-- pipeline/pipeline.conf     # inputs wired, filter{} empty: your canvas
|-- scripts/check.sh           # the referee
`-- LAB2.pdf                   # this document
\end{lstlisting}

\begin{lstlisting}[language=bash]
docker compose up -d
python3 generator/generate_logs.py        # keep it running in a terminal
docker compose logs -f logstash           # watch your rubydebug output
\end{lstlisting}

Three sources land in \code{logs/}, already mounted into the Logstash
container at \code{/var/log/shopfast/}:

\begin{lstlisting}
access.log : 81.2.69.142 - - [26/Jul/2026:01:47:03 +0000] "GET /checkout HTTP/1.1" 502 1370 "-" "Mozilla/5.0..."
app.log    : 2026-07-26T01:47:11Z|payment|ERROR|order=48221|amount=59.90|msg=card declined
orders.csv : 48222,2026-07-26T01:48:02Z,FR,3,129.50
\end{lstlisting}

\begin{trapbox}
About 2\% of the lines in \emph{every} source are malformed. That is not a bug
in the generator — it is the whole point of missions 6 and 7. Plan for it from
the start.
\end{trapbox}

\mission{First flow, and a mystery}{25 min}

\begin{enumerate}
  \item Bring the stack up, start the generator, and confirm raw events from
        \textbf{all three sources} appear in the Logstash output.
  \item Now run \code{docker compose restart logstash}. Observe what is re-read
        and what is not. Then \code{docker compose down \&\& docker compose up -d}.
        Observe again.
  \item \textbf{Write-up}: explain both observations with the correct Logstash
        concept, and name the setting you would use to force a full replay
        during development.
\end{enumerate}

\begin{whybox}[Expected result]
Events from the three files visible in \code{docker compose logs -f logstash};
a correct explanation of what survives a restart, what survives a
\code{down}, and why.
\end{whybox}

\mission{Tame the nginx source}{45 min}

For events coming from \code{access.log} (tag \code{web}), produce documents
where:
\begin{itemize}
  \item the line is fully parsed — no giant \code{message}-only events;
  \item \code{status} and \code{bytes} are \textbf{integers}, named exactly like that;
  \item \code{@timestamp} is the time \textbf{inside the log line}, not the ingestion time;
  \item the raw parsed timestamp string no longer pollutes the document.
\end{itemize}

\begin{whybox}[Expected result]
Typed fields and a correct \code{@timestamp} in rubydebug. The two \emph{M2}
lines of \code{check.sh} turn green once mission 6 ships to Elasticsearch.
\end{whybox}

\begin{docbox}[Doc pointers]
Logstash reference: \emph{grok filter}, \emph{grok patterns list},
\emph{date filter}, \emph{mutate filter}.
\end{docbox}

\mission{Tame the app source}{35 min}

For events from \code{app.log} (tag \code{app}), produce documents with:
\begin{itemize}
  \item fields \code{service}, \code{level}, \code{order} (integer),
        \code{amount} (float), \code{msg};
  \item \code{@timestamp} taken from the line;
  \item parsing implemented with \textbf{grok \emph{or} dissect — your choice}.
\end{itemize}

\textbf{Write-up}: justify the choice in 3 lines (think: format regularity,
failure behavior, performance). A ``both would work, here is why I picked X''
answer is a good answer.

\mission{Tame the CSV source}{25 min}

For events from \code{orders.csv} (tag \code{orders}):
\begin{itemize}
  \item columns are \code{order\_id}, \code{ts}, \code{customer\_country},
        \code{items}, \code{amount};
  \item \code{items} is an integer, \code{amount} a \textbf{float};
  \item \code{@timestamp} comes from the \code{ts} column.
\end{itemize}

\begin{whybox}[Expected result]
Typed order documents in rubydebug; the \emph{M4} line turns green later.
\end{whybox}

\begin{docbox}[Doc pointers]
Logstash reference: \emph{csv filter}.
\end{docbox}

\mission{Enrich and clean}{30 min}

\begin{enumerate}
  \item Enrich web events with \textbf{geolocation} derived from the client IP
        (a \code{geoip} object with at least country and coordinates).
  \item The \code{/healthz} probes are internal monitoring noise: they must
        \textbf{not exist} in your output. Anywhere.
  \item Bonus polish: parse the browser/OS out of the user-agent string.
\end{enumerate}

\begin{whybox}[Expected result]
Web events carry \code{geoip.country\_name} and coordinates; zero
\code{/healthz} documents. \code{check.sh} \emph{M5} lines green later.
\end{whybox}

\mission{Ship it. Lose nothing.}{40 min}

Replace/extend the output section so that:
\begin{itemize}
  \item each source goes to its own dated index:
        \code{shopfast-web-YYYY.MM.dd}, \code{shopfast-app-YYYY.MM.dd},
        \code{shopfast-orders-YYYY.MM.dd};
  \item \textbf{any} event that failed parsing — from any source — lands in
        \code{shopfast-deadletter};
  \item no event path leads nowhere: \textbf{nothing may be silently lost}.
\end{itemize}

\begin{lstlisting}[language=bash]
./scripts/check.sh        # the goal: every line PASS
\end{lstlisting}

\begin{trapbox}
Date-pattern index names use the \emph{event's} \code{@timestamp}. If missions
2--4 got the date wrong, your indices will be named after the wrong day —
check \code{GET \_cat/indices?v} with Course 1 eyes.
\end{trapbox}

\mission{Prove it}{15 min}

\begin{enumerate}
  \item Let the generator run a few minutes, then run \code{check.sh}:
        \textbf{all PASS}.
  \item In Dev Tools, write one query that counts dead-letter events \emph{per
        source file} — paste it and its result into your write-up.
  \item Skim 3 dead-letter documents: can you tell \emph{why} each one failed?
        One line each.
\end{enumerate}

\begin{whybox}[Why this matters]
In the final project, dead-letter handling and ``nothing silently lost'' are
worth real points on an \emph{unseen} dataset — this mission is your rehearsal.
\end{whybox}

\section*{Bonus — One pipeline per source (+2)}

Split your monolithic \code{pipeline.conf} into three isolated pipelines using
\code{pipelines.yml}, without changing any behavior (\code{check.sh} must
still be all green).

\emph{Hint-free zone. The documentation page you need is about multiple
pipelines — and remember how the Docker image decides what to run.}

\bigskip
\begin{center}
\emph{Ship \code{pipeline.conf} + write-up before the deadline.\\
Next session: Course 3 — Kibana, where your clean indices finally become pictures.}
\end{center}

\end{document}
